\section{Projective representations}\label{Section 2}

In this section we recall some basic definitions and properties of projective representations that
will be used throughout this article.

\subsection{Basic definitions}

\begin{Def}
Let $G$ be a finite group and $V$ a finite dimensional complex
vector space. A~map $\rho\colon G\rightarrow {\rm GL}(V)$ is called a
\textit{projective representation} of $G$ if there exists a
function
\begin{gather*}
\alpha\colon \ G\times G \rightarrow \C^{\ast}
\end{gather*}
such that
\begin{gather}
\rho(g)\rho(h)=\alpha(g,h)\rho(gh)\label{ecu:rep-proj}
\end{gather}
for all $g,h\in G$ and $\rho(1)={\rm Id}_{V}$.
\end{Def}
	
Note that if $\rho$ satisfies equation~(\ref{ecu:rep-proj}) then the function $\alpha$ defines a
$\C^{\ast}$-valued normalized 2-cocycle on $G$; that is, for all $g,h,k\in G$ we have:
\begin{gather*}
\alpha(gh,k)\alpha(g,h)=\alpha(g,hk)\alpha(h,k),
\\
\alpha(g,1)=\alpha(1,g)=1.
\end{gather*}
To stress the dependence of $\rho$ on $V$ and $\alpha$, we shall
often refer to $\rho$ as an $\alpha$-representation of $G$
on the space $V$ or, simply as an $\alpha$-representation of
$G$, if $V$ is not pertinent to the discussion.

\begin{Rem}
If $\alpha$ is the trivial cocycle; that is, if $\alpha(g,h)=1$ for all $g,h\in G$, then
$\alpha$-representations of $G$ are simply ordinary representations of $G$.
\end{Rem}

\begin{Def}
Suppose that $\rho_{i}\colon G\rightarrow {\rm GL}(V_{i})$ with $(i=1,2)$ are two $\alpha$-representations.
\begin{enumerate}\itemsep=0pt
\item A linear map $\varphi\colon V_{1}\to V_{2}$ is said to be a map of projective representations or a
$G$-map if~for any $g\in G$ and any $v\in V_{1}$ we have $\varphi(\rho_{1}(g)v)=\rho_{2}(g)\varphi(v)$.
We write $\Hom_{G}(V_{1},V_{2})$ for the set of $G$-morphisms from $V_{1}$ to $V_{2}$.

\item The $\alpha$-representations $V_{1}$ and $V_{2}$
are said to be \textit{linearly equivalent} or \textit{isomorphic} if
there exists a $G$-map $f\colon V_{1}\to V_{2}$ that is an isomorphism of vector spaces.
In other words, $V_{1}$ and $V_{2}$ are isomorphic if there is a vector space
isomorphism $f\colon V_{1}\rightarrow V_{2}$ such that $\rho_{2}(g)=f\rho_{1}(g)f^{-1}$ for all $g\in G$.
\end{enumerate}
\end{Def}

Just as in the case of the ordinary theory of representations of groups
we have similar notions for projective representations such as irreducible representations
and unitary representations. Given a $\C^{\ast}$-valued normalized 2-cocycle $\alpha$ on
$G$ we denote by $\Irr^{\alpha}(G)$ the set of isomorphism classes of
complex irreducible $\alpha$-representations of $G$.
If $\rho\colon G\to {\rm GL}(V)$ is an irreducible $\alpha$-representation of $G$ then $[\rho]\in \Irr^{\alpha}(G)$
denotes the corresponding isomorphism class. Observe that the direct sum of two $\alpha$-representations
is also an $\alpha$-representation. Thus we can form the monoid of isomorphism classes of $\alpha$-representations.
The $\alpha$-twisted representation group of $G$, denoted by $R^{\alpha}(G)$,
is defined as the associated Grothendieck group. As an abelian group $R^{\alpha}(G)$ is a free abelian
group with one generator for each element in $\Irr^{\alpha}(G)$.
We remark that the classical results of representation theory
such as complete reducibility and Schur's lemma also hold for the case of projective representations.
We refer the reader to~\cite{Karpilovsky2} for
the basic theory of~projective representations.

